\documentclass[12pt]{article}
\usepackage[T1]{fontenc}
\usepackage[utf8]{inputenc}
\usepackage{lmodern}
\usepackage{textcomp}
\usepackage{enumitem}
\usepackage{geometry}
\geometry{margin=1in}
\begin{document}
\begin{center}
Answer Key - Mathematics - K-12 Level Assessment 3 \
\small (Answers are ordered exactly as in the question paper.)
\end{center}

\section*{Multiple Choice}
\begin{enumerate}[label=\arabic*.]
\item \textbf{Answer:} C
\\ \textit{Options:} A) 16; B) 20; C) 24; D) 28
\\ \textit{Note:} Correct option: C - 24
\item \textbf{Answer:} A
\\ \textit{Options:} A) 4.5; B) 2.25; C) 6.6; D) 3.3
\\ \textit{Note:} Correct option: A - 4.5
\item \textbf{Answer:} C
\\ \textit{Options:} A) 100; B) 999; C) 625; D) 500
\\ \textit{Note:} Correct option: C - 625
\item \textbf{Answer:} B
\\ \textit{Options:} A) central tendency.; B) variability.; C) symmetry.; D) skewness.
\\ \textit{Note:} Correct option: B - variability.
\item \textbf{Answer:} D
\\ \textit{Options:} A) 5; B) 11; C) 17; D) 199
\\ \textit{Note:} Correct option: D - 199
\end{enumerate}

\section*{True / False}
\begin{enumerate}[label=\arabic*.]
\setcounter{enumi}{5}
\item \textbf{Answer:} False
\\ \textit{Options:} A) True; B) False
\\ \textit{Note:} LLM-generated false statement. Correct answer: '1 and 1 over 2 ft'.
\item \textbf{Answer:} True
\\ \textit{Options:} A) True; B) False
\\ \textit{Note:} LLM-generated true statement based on MCQ.
\item \textbf{Answer:} False
\\ \textit{Options:} A) True; B) False
\\ \textit{Note:} LLM-generated false statement. Correct answer: '7'.
\item \textbf{Answer:} True
\\ \textit{Options:} A) True; B) False
\\ \textit{Note:} LLM-generated true statement based on MCQ.
\item \textbf{Answer:} False
\\ \textit{Options:} A) True; B) False
\\ \textit{Note:} LLM-generated false statement. Correct answer: '5.7'.
\end{enumerate}

\section*{Long Form Response}
\begin{enumerate}[label=\arabic*.]
\setcounter{enumi}{10}
\item \textbf{Answer:} Note that \[a_k = \frac{1}{k^2+k} = \frac{1}{k} - \frac{1}{k+1}\]for each $k.$ Thus, the sum telescopes: \[\begin{aligned} a_m + a_{m+1} + \dots + a_{n-1} & = \left(\frac{1}{m} - \frac{1}{m+1}\right) + \left(\frac{1}{m+1} - \frac{1}{m+2}\right) + \dots + \left(\frac{1}{n-1} - \frac{1}{n}\right) \\ &= \frac{1}{m} - \frac{1}{n}. \end{aligned}\]Therefore, we have the equation $1/m - 1/n = 1/29.$ Multiplying by $29 mn$ on both sides, we have $29 n - 29 m = mn,$ or $mn + 29 m - 29 n = 0.$ Subtracting $29^2$ from both sides, we get \[(m-29)(n+29) = -29^2.\]Since $29$ is prime and $0 < m < n,$ the only possibility is $m-29 = -1$ and $n+29 = 841,$ which gives $m = 28$ and $n = 812.$ Thus, $m+n=28+812=\boxed{840}.$
\\ \textit{Note:} Students should show all steps in their mathematical solution.
\item \textbf{Answer:} 2 \ensuremath{\leq} y \ensuremath{\leq} 8. Explain why this is the correct answer and provide supporting evidence. Show the calculation or reasoning that leads to this conclusion.
\\ \textit{Note:} Long-form derived from MCQ; award credit for stating the correct idea and giving supporting reasoning.
\end{enumerate}

\vfill
\noindent\textit{End of Answer Key}
\end{document}